\section*{Part I: What a Physical Asset Is Supposed to Do}
\addcontentsline{toc}{section}{Part I: What a Physical Asset Is Supposed to Do}


\subsection{1. The mistake of beginning with an acreage target}

A campus can be a laboratory of human relationships and scientific work, but its value is not measured by acres. For a new technical institution, possible outputs include hands-on experiments, peer interaction, professional identity, faculty collaboration, student housing, research safety, demonstration to institutional buyers, and access to networks. Each requires different physical conditions. A single research facility close to collaborators can outperform a large remote campus; a studio may be sufficient for a small program whose learning happens online.

Write the \textbf{physical mission statement}: which activity cannot be provided at comparable quality through rented rooms, partners, online tools, or short residencies? If the answer is vague, postpone irreversible commitments without abandoning physical ambition.


\subsection{2. 6 different assets hiding under the word campus}

A founder may need: (1) periodic classroom access; (2) a branded teaching studio; (3) specialized shared laboratory facilities; (4) academic offices and faculty collaboration space; (5) residential student and faculty accommodation; or (6) a legally recognized higher-education campus meeting statutory conditions. These have different cost curves, regulatory pathways, and geographic constraints. A location excellent for executive workshops in central Mumbai may be unsuitable for an expansive residential research campus. A cheap remote parcel may be attractive for a warehouse and terrible for attracting scientists.

Separate these functions in the property brief before searching. Do not let the existence of land dictate the institution's operating model.


\subsection{3. Physical expansion is a compound option}

A site purchase creates exposure to land price movements, local law, infrastructure changes, finance costs, and development approvals. It can also preserve options for future buildings. An option is attractive only if the cost of holding it is proportionate to the future flexibility it creates. Purchasing an illiquid plot with high carrying costs may reduce, rather than increase, the founder's options.

An early institute can acquire \textit{use rights} before acquiring \textit{ownership rights}: short-term bookings, partner laboratory access, a serviced office, a conditional lease, a build-to-suit arrangement, or a purchase option with legal safeguards. The real comparison is among the states of the world the firm can survive.


\subsection{4. 3 time horizons}

\textbf{Near term:} proof of demand and specific physical need. \textbf{Medium term:} predictable recurring utilization, repeatable operations and a stable management team. \textbf{Long term:} research and degree-related physical commitments, permanent faculty communities, and asset stewardship. A 30-year facility decision should not be made using only next year's enrollment estimate, but nor should an imagined Year-30 university justify making the Year-1 enterprise insolvent.

Model each horizon separately. Costs may be committed now while benefits emerge much later. That is precisely why property can consume optionality.


\subsection{5. Ownership, leasing, partnering and managed access}

Ownership provides the greatest nominal control but exposes the firm to acquisition, maintenance, compliance, and illiquidity. Leasing may preserve capital yet impose escalation, renewal, restoration, access, and landlord risks. Partnership can provide specialized labs but creates dependence on another institution. Pay-per-use maximizes flexibility while limiting permanent identity and assured scheduling. Build-to-suit transfers some construction responsibility while creating long-term rental commitments.

The cheapest headline rent is not necessarily the cheapest total option. A properly designed comparison includes economic cost, service levels, control, rights, future capacity, reversibility, and legal feasibility.


\subsection{6. The location strategy should precede the site search}

A founder must make 3 linked decisions: \textbf{region or city}, \textbf{micro-market}, and \textbf{parcel}. Region determines research/employer ecosystem and overall regulatory context; micro-market determines accessibility, housing and transport; parcel determines legal risk, frontage, topography, utilities, buildability and acquisition cost. Scoring a parcel before choosing the right region is backward reasoning.

A physical campus's benefits depend partly on proximity to complementary assets that the institute cannot build: hospitals, schools, airports, employers, other universities, skilled vendors, entertainment, public transport and talent networks.
